%========================= PARTE B: categorías 6–10 ==========================

\section{Métodos numéricos}

\Q{nivel básico · lámina 12}{¿Por qué volúmenes finitos y no elementos
finitos?}
\A{Por el tipo de solución. FEM, con su formulación débil global, es óptimo
para dominios complejos con soluciones suaves. En flujos relativistas
emergen choques que vuelven singular la forma diferencial: se necesitan
soluciones débiles con saltos, y la formulación integral de FVM las admite
por construcción; los flujos numéricos de Riemann respetan las condiciones
de salto de Rankine--Hugoniot (Toro 2009). Además la conservación discreta
es exacta: el flujo que sale de una celda entra en la vecina, requisito
para capturar choques con las velocidades correctas.}

\Q{nivel medio · lámina 13}{Enuncie el teorema de Godunov y explique cómo lo
esquiva MP5.}
\A{Ningún esquema lineal de orden mayor que uno preserva monotonicidad:
alto orden lineal implica oscilaciones de Gibbs en discontinuidades. La
salida es la no linealidad del esquema: limitadores. Los TVD de segundo
orden (minmod y afines) son monótonos pero recortan los extremos suaves
---y con ello destruyen la vorticidad que queremos medir---. MP5 (Suresh \&
Huynh 1997) usa un \emph{stencil} de cinco puntos con límite de
preservación de monotonicidad mediante una mediana proyectiva: mantiene
quinto orden en regiones suaves, no oscila en choques y retiene los
extremos de vorticidad de forma casi espectral.}

\Q{nivel básico · lámina 13}{¿Qué es el flujo HLL y cuáles son sus
propiedades?}
\A{Un solucionador de Riemann aproximado de dos ondas: supone un único
estado intermedio entre las velocidades características extremas
$\lambda_L,\lambda_R$ y da el flujo
$\mathbf F^{hll}=[\lambda_R\mathbf F_L-\lambda_L\mathbf F_R+
\lambda_L\lambda_R(\mathbf U_R-\mathbf U_L)]/(\lambda_R-\lambda_L)$
(Harten, Lax \& van Leer 1983). Es robusto, positivo, libre de jacobianos
---valioso con un sistema de 14 variables--- y entrópico; su precio es
difusión sobre la onda de contacto, que en nuestro esquema compensa la
agudeza de los estados reconstruidos a quinto orden.}

\Q{nivel alto · respaldo B7}{El artículo del código presenta un solucionador
HLLC. ¿Por qué el barrido usa HLL?}
\A{Decisión de diseño documentada. HLLC restaura la onda de contacto
explícitamente; pero Miranda-Aranguren et al.\ (2018) muestran que HLL
alimentado con estados MP5 alcanza una nitidez de contacto comparable: la
información que HLLC añade en el solucionador, MP5 la aporta en la
reconstrucción. Para un barrido de 62 simulaciones elegimos la combinación
más robusta (positividad garantizada, sin estados intermedios delicados) y
más barata por paso. La verificación de que la difusión numérica resultante
quedó por debajo de la física es empírica: el plateau de $J_{\max}$.}

\Q{nivel alto · láminas 13 y 26}{¿Cómo garantiza que su vórtice es física y
no viscosidad artificial?}
\A{En tres niveles. Diseño: MP5$+$HLL se eligió para que la disipación
numérica quede por debajo de la física en todo el barrido. Evidencia
empírica: el plateau de $J_{\max}$ localiza dónde la difusividad física
$\eta$ supera el error de truncamiento; la zona de transición completa está
en ese régimen. Conceptual: la KHI 2D \emph{ideal} ni siquiera converge con
la resolución (Lecoanet 2016), de modo que es precisamente la $\eta$
explícita la que da significado convergente y bien planteado a
$\gamma_{\rm KHI}(\sigma)$.}

\Q{nivel medio · lámina 14}{Describa la partición IMEX y el costo real de su
parte implícita.}
\A{Los flujos hiperbólicos se integran explícitamente (MP5$+$HLL, $\Delta t
\propto\Delta x$); la fuente rígida $-\sigma\mathbf E$, implícitamente, en
un Runge--Kutta IMEX de tipo SSP (Pareschi \& Russo 2005). Como la fuente
no acopla celdas vecinas, la etapa implícita es un sistema algebraico local
$3\times3$ para $\mathbf E$ con solución analítica: no hay Newton global ni
inversión iterativa, y el sobrecosto por paso es despreciable frente al
salto de paso temporal que se gana ---de $\Delta t\propto\Delta x^2/\sigma$
a $\Delta t\propto\Delta x$.}

\Q{nivel medio · lámina 15}{¿Cómo recupera el código las variables
primitivas y qué guardias tiene?}
\A{El mapeo conservadas$\to$primitivas no tiene inversa cerrada porque $W$
acopla cinemática y termodinámica. Cueva lo reduce, con la EoS politrópica,
a una cuártica en $W$ que se resuelve analíticamente (Tschirnhaus $+$
Ferrari--Cardano) y se refina con pocas iteraciones de Newton--Raphson.
Guardias: rechazo de $v^2\ge1$ y de $p\le0$, rama linealizada estable a
bajas velocidades y \emph{fallback} a promedios de celdas vecinas si la
raíz es inadmisible. Se ejecuta en cada celda y subpaso: su robustez define
el límite práctico del barrido.}

\Q{nivel medio · lámina 17}{¿Qué recursos computacionales consumió el
trabajo y en qué arquitectura?}
\A{Unas $3\times10^4$ horas-núcleo acumuladas en CPU. Cueva es Fortran
paralelo sobre CPU ---no CUDA---, con volúmenes finitos unigrid; el
análisis posterior es Python. La elección es deliberada: para mallas
$512\times256$ en 2D el cuello es el barrido de parámetros, no la
simulación individual, y la robustez portable de CPU simplifica reproducir
cualquier punto del barrido.}

\Q{nivel alto · lámina 19}{¿Hicieron estudio de convergencia de malla? ¿Qué
responde a la exigencia de independencia de malla?}
\A{Con matiz: para la KHI 2D ideal el estudio de convergencia clásico está
condenado ---Lecoanet et al.\ (2016) demuestran que sin disipación
explícita las cantidades no convergen al refinar---. Nuestro planteamiento
invierte la lógica: la $\eta$ física fija la escala disipativa, y la
evidencia de resolución suficiente es que la escala resistiva esté resuelta
donde extraemos conclusiones cuantitativas: el plateau de $J_{\max}$ y las
$\approx13$ celdas por capa de cizalla lo acreditan. Declaramos que a las
$\sigma$ más altas las láminas rozan la malla, y por eso el exponente
$1.22$ se presenta con esa salvedad.}

\section{Montaje experimental y diseño de campañas}

\Q{nivel básico · lámina 16}{Describa su condición inicial. ¿Qué gatilla la
inestabilidad?}
\A{Doble capa de cizalla tipo Mizuno (2013): perfiles
$V_y(x)=\pm v_{\rm sh}\tanh[(x\mp0.5)/a_{kh}]$ con $v_{\rm sh}=0.5$ y
$a_{kh}=0.05$; contraste de densidad $\eta_\rho=0.1$ (chorro diluido en
ambiente denso); campo $\mathbf B=(0,\sqrt{0.02},\sqrt{2.0})$;
$\mathbf E(0)=0$. La semilla es una perturbación senoidal en $V_x$, el modo
fundamental $k=2\pi/L_y$ con confinamiento gaussiano en las capas. El
contraste de densidad modela la configuración chorro--medio pero no gatilla
nada: sin la perturbación de velocidad el estado base es estacionario.}

\Q{nivel medio · lámina 16}{¿Por qué eligieron un campo guía $B_z$
dominante?}
\A{Es la decisión topológica que hace posible el experimento: con
$\mathbf k$ en $\hat y$ y el campo casi todo en $\hat z$, el producto
$\mathbf k\cdot\mathbf B\propto B_y$ es débil, de modo que el campo aporta
presión magnética pero casi ninguna tensión. Un campo fuerte paralelo a
$\mathbf k$ habría estabilizado la KHI y no habría nada que medir. Así la
inestabilidad se desarrolla y el efecto de $\sigma$ queda aislado; la
componente $B_y=\sqrt{0.02}$ pequeña conserva la semilla de dínamo que
señala Mizuno.}

\Q{nivel medio · lámina 16 · respaldo B6}{¿Las dos capas de cizalla no se
contaminan mutuamente? ¿Y las fronteras?}
\A{No: el acople entre capas decae como $e^{-kD}$ con la separación $D$;
para nuestra geometría es del orden del $0.2\%$, verificado con el
solucionador lineal (respaldo B6). Las condiciones de frontera son
periódicas en $y$ ---compatibles con el modo fundamental--- y abiertas en
$x$, con el dominio suficientemente ancho para que las ondas emitidas
salgan sin reflejarse sobre la capa durante el horizonte simulado.}

\Q{nivel básico · lámina 16}{¿Por qué $\Gamma=4/3$?}
\A{Es el índice adiabático del gas relativista caliente (fotones o pares, o
bariones con presión dominada por especies relativistas), el régimen
apropiado para el interior de chorros de AGN. La EoS es politrópica por
diseño: Mizuno (2013) muestra que la elección de EoS modula la KHI
resistiva, así que fijarla elimina un eje de variación y deja $\sigma$ y
$\mathbf B$ como únicos parámetros de control.}

\Q{nivel básico · lámina 17}{Presente su matriz de simulaciones.}
\A{Campaña 1: barrido de 44 conductividades $\sigma\in[100,10\,500]$ con
CFL $0.1$, todo lo demás fijo. Campaña 2: a $\sigma=6000$ (CFL $0.04$) y
$\sigma=10\,000$ (CFL $0.02$), tres familias: A ---intensidad $B_0=1.0,
1.5,2.0$---, B ---rotación del campo en el plano $y$--$z$, que activa
tensión--- y C ---rotación en $x$--$z$ con $\mathbf k\cdot\mathbf B=0$, sin
tensión---. Todas con malla $512\times256$ y horizonte $t>5$.}

\Q{nivel medio · lámina 17}{¿Qué datos excluyeron y con qué criterio?}
\A{Tres exclusiones, declaradas: en la campaña A, $B_0=0.25$ y $0.5$
fallaron por inestabilidad numérica antes de la fase no lineal; en la B,
$\theta=15^\circ$ se detuvo por fallo de convergencia del solucionador de
primitivas; y en el barrido, el análisis cuantitativo de $\gamma$ usa
$\sigma\ge1600$ con $R^2\ge0.90$ y señal-ruido $\ge3$ ---por debajo no
existe fase lineal que ajustar, no es un ajuste pobre---. Los casos
excluidos se conservan y se reportan: la transparencia es parte del
resultado.}

\Q{nivel medio · lámina 16}{¿Es consistente iniciar con $\mathbf E=0$ si la
ley de Ohm exige otro campo?}
\A{Sí, y es deliberado: el término rígido $-\sigma(\mathbf E-
\mathbf E_{\rm ohm})$ relaja el campo eléctrico hacia el valor consistente
con la conductividad en un tiempo $\sim1/\sigma$, órdenes de magnitud más
corto que el crecimiento de la KHI. El transitorio inicial queda fuera de
la ventana de ajuste $t\in[2.4,3.4]$ precisamente para que esta relajación
---y el acomodo acústico inicial--- no contaminen la medición.}

\Q{nivel alto · lámina 17}{¿Por qué el rango $[100,10\,500]$ y no
conductividades mayores?}
\A{El extremo inferior garantiza cubrir el régimen dominado por difusión; el
superior está dictado por la resolución: a $\sigma$ mayores la capa
resistiva $\delta\sim S^{-1/2}$ cae por debajo de la malla y los
resultados dejarían de estar en el régimen controlado que exigimos
(disipación física $>$ numérica). Con $512\times256$, $\sigma=10\,500$ es
la frontera prudente; medir el plateau ideal directamente exige
$\sigma\sim2$--$3\times10^4$ con más resolución, y así se declara como
trabajo futuro.}

\section{Observables y metodología de análisis}

\Q{nivel básico · lámina 11}{Defina su observable primario.}
\A{La enstrofía de perturbación. A la vorticidad $\omega_z=\partial_xV_y-
\partial_yV_x$ se le sustrae el promedio en $y$ de la configuración
inicial, para remover la cizalla base: $\omega_z'=\omega_z-\langle\omega_z
(t{=}0)\rangle_y$. Se integra su cuadrado: $\Omega_{zp}=\int(\omega_z')^2
dA$. Mide la intensidad rotacional de la perturbación, es positiva,
integral (robusta al ruido puntual) y estándar en la literatura de KHI
resistiva (Nastro 2022).}

\Q{nivel básico · lámina 11}{¿Por qué la pendiente de $\ln\Omega_{zp}$ es
$2\gamma$ y no $\gamma$?}
\A{Porque en fase lineal cada modo crece como $e^{\gamma t}$ y la enstrofía
es \emph{cuadrática} en la amplitud de la perturbación: su logaritmo crece
con el doble de la tasa. Hacerlo explícito evita el error de factor 2 más
común al comparar tasas entre trabajos que usan observables lineales
(amplitud de modo) y cuadráticos (energías, enstrofías).}

\Q{nivel medio · lámina 20}{¿Cómo eligieron la ventana de ajuste
$t\in[2.4,3.4]$?}
\A{Es posterior al transitorio inicial ---relajación de Ohm y acomodo
acústico de la condición inicial--- y anterior a la saturación del modo
fundamental incluso en los casos más rápidos. Se usa la \emph{misma}
ventana para las 44 simulaciones: una ventana fija sacrifica algo de ajuste
en los casos extremos, pero garantiza que las tasas sean comparables entre
sí y no un artefacto de ventanas elegidas a conveniencia.}

\Q{nivel básico · lámina 11}{Enumere sus diagnósticos complementarios y qué
mide cada uno.}
\A{$\Omega_{\rm tot}$ y la fracción fuera del plano $f_{Vz}=(\Omega_{\rm
tot}-\Omega_z)/\Omega_{\rm tot}$: estructuras secundarias con velocidad
$V_z$. $J_{\max}$: el adelgazamiento de las láminas de corriente.
$\int\mathbf E\cdot\mathbf J\,dA$: disipación óhmica global. Las energías
$E_{\rm kin}$, $E_{\rm mag}$, $E_{\rm int}$ y su suma: canales de
conversión y control de conservación. Y la tasa normalizada
$\hat\omega=\gamma\,a_{kh}/v_{\rm sh}$ para comparar con la literatura.}

\Q{nivel medio · lámina 11}{¿Por qué enstrofía y no la amplitud de un modo de
Fourier?}
\A{La amplitud modal es óptima para un problema estrictamente lineal, pero
aquí el fondo se difunde (el perfil $\tanh$ se ensancha con $\eta$), lo que
contamina los modos bajos, y en el régimen resistivo la energía se
redistribuye entre armónicos desde tiempos tempranos. La enstrofía de
perturbación integra toda la actividad rotacional de la perturbación, es
insensible a esa redistribución y sigue siendo válida cuando la fase lineal
es marginal ---exactamente el régimen que queríamos clasificar.}

\Q{nivel medio · lámina 26}{¿Qué significa físicamente el máximo de
$f_{Vz}$ en la transición?}
\A{$f_{Vz}$ mide cuánta enstrofía vive en movimientos fuera del plano. En el
régimen resistivo la difusión suprime los gradientes que alimentarían
$V_z$; en el cuasi-ideal el congelamiento rigidiza el plasma y también las
suprime. El máximo ($f_{Vz}=0.37$ en $\sigma\approx1400$) aparece donde
difusión y advección compiten: la zona de transición es el hábitat de las
estructuras secundarias. Es uno de los cuatro relojes que usamos para
delimitar la zona crítica.}

\section{Resultados del barrido: fase lineal y transición}

\Q{nivel básico · lámina 21}{Presente su resultado paramétrico central.}
\A{La tasa lineal en función de la conductividad queda descrita por un
sigmoide con desplazamiento:
$\gamma_{\rm KHI}(\sigma)=C+\gamma_0/(1+e^{-k(\sigma-\sigma_0)})$. La
asíntota ideal vale $C+\gamma_0=1.03\pm0.05$ (normalizada,
$\hat\omega=0.103$) y el centro de la transición es
$\sigma_0=2815\pm123$. El modelo se ajusta sobre las 35 simulaciones con
fase lineal medible y supera decisivamente al sigmoide sin desplazamiento.}

\Q{nivel medio · lámina 21}{El desplazamiento $C$ del ajuste es negativo.
¿Una tasa de crecimiento negativa?}
\A{No: $C$ aislado no es un observable físico, es un parámetro del modelo
---la extrapolación formal del sigmoide hacia $\sigma\to-\infty$, fuera del
dominio físico---. Lo defendible es su papel funcional: con $C$ libre el
modelo predice un 36\% mejor los nueve puntos resistivos excluidos del
ajuste. Además $C$ y $\gamma_0$ están fuertemente anticorrelacionados
($r=-0.997$); la combinación robusta es la suma $C+\gamma_0$, que varía
menos del 4\% entre variantes del ajuste y es la que interpretamos.}

\Q{nivel medio · lámina 22}{¿Cómo estimaron $\sigma_{\rm crit}$ y por qué
cuatro veces?}
\A{Con un método común aplicado a cuatro observables independientes: el
punto de máxima derivada respecto de $\log_{10}\sigma$ (inflexión), con
error el máximo entre el remuestreo jackknife y el semiancho de la zona de
inflexión. Resultados: $2815\pm123$ (sigmoide de $\gamma$), $3400\pm2175$
($\Omega^{\max}_{zp}$), $6000\pm1147$ ($t_{\rm peak}$) y $6800\pm2625$
($\gamma$ en escala log). Cuatro relojes distintos porque cada fenómeno se
``idealiza'' a su propia conductividad: esa es la física del resultado.}

\Q{nivel básico · lámina 23}{¿Por qué insisten en que $\sigma_{\rm crit}$ es
una zona y no un número?}
\A{Porque la dispersión entre estimadores ($\pm1944$) excede en un orden de
magnitud el error formal del mejor ajuste ($\pm123$): reportar
``$2815\pm123$'' sería precisión espuria. La lectura honesta es una zona de
transición $\sigma\in[1400,7000]$ ---del máximo de $f_{Vz}$ a la
saturación no lineal--- con centro $\approx2815$. La dispersión no es
ruido: es la firma de una transición multiescala en la que cada observable
se congela a su escala. En número de Lundquist, $S\sim10^3$.}

\Q{nivel medio · lámina 20}{¿Cómo clasificaron los regímenes del barrido?}
\A{Por la calidad del ajuste lineal en la ventana canónica: $R^2<0.90$
define el régimen resistivo (no hay recta que ajustar), $0.90$--$0.995$ la
transición y $\ge0.995$ el cuasi-ideal. Es un criterio operativo y
reproducible que coincide con la morfología: difusión de vórtices en el
resistivo, subestructura fina y turbulencia de pequeña escala en el
cuasi-ideal.}

\Q{nivel medio · lámina 20}{Para $\sigma\lesssim1400$ no reportan tasa.
¿Fracasó el método?}
\A{No: el resultado es que \emph{no existe} fase lineal que medir. Con
difusividad grande, la capa se ensancha en el tiempo de crecimiento del
modo: la difusión domina desde $t=0$ y el crecimiento exponencial nunca se
establece. Forzar una pendiente daría un número sin significado. Esas
simulaciones se conservan para la clasificación de regímenes y para los
diagnósticos globales, que no requieren fase lineal.}

\Q{nivel básico · lámina 19}{Describa las tres fases temporales que exhibe
el barrido.}
\A{Fase lineal hacia $t\sim3$: la perturbación crece exponencialmente y la
interfaz se ondula. Enrollamiento no lineal hacia $t\sim8$: los vórtices
maduran, generan láminas de corriente y, según $\sigma$, reconectan o
enrollan el campo. Saturación hacia $t\sim14$: la enstrofía alcanza su
máximo y decae mientras la capa se ensancha. Las 44 simulaciones muestran
la secuencia; lo que cambia con $\sigma$ es la tasa, la morfología y el
destino de la energía.}

\Q{nivel alto · lámina 21}{¿Qué garantiza que el sigmoide no es un
sobreajuste con 4 parámetros?}
\A{Tres controles. Comparación anidada: contra el sigmoide de 3 parámetros
($C=0$), sobre los mismos datos, $\Delta\mathrm{AIC}=-126$; en la escala de
Burnham \& Anderson, evidencia decisiva que paga con creces el parámetro
extra. Validación fuera de muestra: mejora del 36\% (RMSE
$0.039\to0.025$) sobre los 9 puntos resistivos no usados. Y honestidad
residual: los residuos están autocorrelacionados, así que no confiamos el
error a la matriz de covarianza del ajuste sino a la dispersión
inter-método de la lámina siguiente.}

\section{Contraste con la teoría lineal}

\Q{nivel básico · lámina 24}{Explique la ``escalera de predicciones''.}
\A{Tres peldaños. Techo hidrodinámico inviscido: la ecuación de Rayleigh
para el perfil $\tanh$ da $\hat\omega_{\max}=0.190$ (Michalke 1964).
Predicción magnetosónica: como $\mathbf k\perp\mathbf B$, el campo aporta
presión pero no tensión y la velocidad de restitución es
$v_{f\perp}=0.691$; la teoría compresible tipo Lees--Lin da entonces
$\hat\omega=0.095$. Valor medido: la asíntota del sigmoide,
$\hat\omega=0.103$. Acuerdo del 8\% sin parámetros de ajuste.}

\Q{nivel medio · lámina 24 · respaldo B6}{¿Cómo validaron su solucionador de
teoría lineal?}
\A{Reproduciendo el resultado clásico: para el perfil $\tanh$ inviscido e
incompresible, el solucionador de Rayleigh recupera
$\hat\omega_{\max}=0.1897$ de Michalke (1964) en el número de onda
correcto. Además verifica el umbral compresible ($\sqrt2\,c_s$ del caso
\emph{vortex sheet}, Landau 1944) y el desacople de las dos capas
($e^{-kD}\approx0.2\%$). Sobre esa base validada se aplica la
sustitución magnetosónica.}

\Q{nivel medio · lámina 24}{Justifique la sustitución $c_s\to v_{f\perp}$.}
\A{Es el resultado de Chow, Rosen \& Piran (2023) para geometría de campo
guía: cuando $\mathbf k\perp\mathbf B$, las perturbaciones comprimen el
campo sin doblarlo; el campo no ejerce tensión restauradora sobre el modo,
pero su presión sí rigidiza el medio. La velocidad efectiva de restitución
pasa de $c_s$ a la magnetosónica rápida perpendicular
$v_{f\perp}=\sqrt{c_s^2+v_A^2}$ relativista, $0.691$ en nuestro estado
base. ``Presión pero no tensión'' resume la física de la sustitución.}

\Q{nivel medio · lámina 24}{¿Cuál es el control negativo de esa cadena?}
\A{Repetir la predicción con $c_s$ en lugar de $v_{f\perp}$: el resultado
colapsa a $\hat\omega\approx0.016$, seis veces por debajo de lo medido.
Eso demuestra que el acuerdo del 8\% no es genérico ---no cualquier
velocidad lo produce--- y que es la presión magnética del campo guía la
que sostiene la inestabilidad en nuestro régimen: sin ella el flujo sería
demasiado supersónico respecto de la restitución puramente acústica.}

\Q{nivel medio · lámina 24}{¿Dónde queda su flujo respecto del umbral de
estabilización relativista?}
\A{El criterio (Königl 1980; Bodo 2004; Chow 2023) es que la discontinuidad
es inestable si el Mach relativista construido con la velocidad de
restitución cumple $0<M_{re}<\sqrt2$. El nuestro es $M_{re}=0.60$:
holgadamente inestable, lejos del umbral $\sqrt2\approx1.41$. Importa
porque la extrapolación de la asíntota no roza ninguna frontera de
estabilidad: el plateau ideal es genuino, no un borde de supresión.}

\Q{nivel alto · lámina 24}{¿Por qué no resolvieron la relación de dispersión
RMHD completa?}
\A{La dispersión compresible magnetizada relativista para el perfil suave
es un polinomio de grado 8 en la frecuencia con coeficientes dependientes
del perfil (Osmanov 2008; Chow 2023): resolverla con nuestros parámetros es
un proyecto en sí mismo, declarado como trabajo futuro. Para el contraste
de este trabajo empleamos la cadena validada Michalke $\to$ sustitución
escalar, y por eso presentamos el resultado como \emph{consistencia
cuantitativa fuerte} ---8\%, sin parámetros libres--- y no como
validación cerrada. La asíntota, además, es extrapolación: el dato máximo
alcanza el 85\% del plateau.}
